\section{Conclusion}\label{sec:conclusion120}
Neural Bellman Operators connect learned continuation operators to feasible
decisions and the economic value of their implemented future behavior. The
construction requires a fixed original recursion and a complete accounting of
the information used by its critic, actor and numerical certificate. The
constrained investment model makes these requirements explicit while retaining
continuous innovations and nonlinear recursive risk.

The new action-comparison theory uses the geometry of the original Bellman
problem. A constrained state-gradient modulus yields an action-scaled
continuation-difference bound; convex integration and a constrained risk-defect
transfer then provide computable comparison intervals. Qualification and a
strict original-policy improvement bound jointly govern replacement of a
completed incumbent. The proof also identifies what additional precision can
resolve and what remains hidden behind a uniform future-policy allowance. This
distinction gives a concrete stopping criterion and a specified information
problem for subsequent NBO algorithms.

The independent experiment certifies gains over the completed learned
incumbents under the original investment objective. All 42 learned-incumbent
refinement groups have positive simultaneous population mean-gain lower bounds;
the independent optimum supports also certify substantial reductions in
original regret and optimal-action error bounds. These gains coexist with a
strong online incumbent that is already close to optimal. The new paired
service changes twelve returned actions relative to inherited safe improvement,
but no cell-level incremental mean advantage is resolved, and all sixteen
primary costs and all forty-eight matched clocks are higher. A mixed comparison
bound changes one recorded acceptance decision. The direct conventional proposal
improves five cell means and is faster than the paired service throughout the
matched clocks. The identical ReLU and hinge policies and the empty ReLU
full-menu adoption regions preserve the unresolved neural contribution as an
explicit scientific problem.

The remaining research priorities concern the joint design of continuation
learning and paid economic decisions. They include acquisition losses that
reduce action-comparison work, information about the actual future controller
that resolves gains below the present allowance, and matched experiments with
longer horizons, multiple controls and higher-dimensional continuous risk.
The retained stochastic-control, recursive-utility and stopped-boundary results
provide separate rigorous foundations for this program. Progress on these
problems must preserve the original objective, give conventional methods the
same usable information, and demonstrate the economic value of the resulting
policies together with their full computational cost.
